\section*{R\'esum\'e}
\label{sec:resume_fr}
\addcontentsline{toc}{section}{R\'esum\'e / Abstract / Resumen}

\begin{otherlanguage}{french}
Dans une cha\^ine d'approvisionnement, l'urgence qu'un op\'erateur d\'eclare diverge souvent
de celle que l'on calcule depuis l'\'etat op\'erationnel r\'eel. Sur-d\'eclarer gaspille une
capacit\'e partag\'ee~; sous-d\'eclarer laisse un risque de rupture invisible. Aucun jumeau
num\'erique, publi\'e ou brevet\'e, ne mesure encore cet \'ecart.

Cette mission, men\'ee \`a la Direction de l'Innovation d'ALTEN, a con\c{c}u et \'eprouv\'e
une couche logicielle qui le mesure, livr\'ee sous le nom de \textbf{SupplyScore}.
L'urgence r\'eelle est agr\'eg\'ee depuis six blocs
d'indicateurs par une r\`egle \emph{non compensatoire}~: le bloc le plus d\'egrad\'e tire le
score vers son maximum, sans qu'un bloc sain puisse l'en emp\^echer. Elle est ensuite
propag\'ee sur le r\'eseau. L'urgence d\'eclar\'ee, elle, est \'elicit\'ee par un
questionnaire de comparaisons par paires contr\^ol\'e en coh\'erence, et l'\'ecart entre les
deux est not\'e sur $[0,100]$.

Dans les tests de \emph{serious game} r\'ealis\'es, sur les six tours pr\'ec\'edant la mise au
jour d'une crise, l'\'ecart d\'esigne trois n\oe{}uds, tous trois futurs d\'efaillants, sans
une seule fausse alerte~; le plus mal not\'e est signal\'e huit tours avant de rater son
jalon, alors que son responsable d\'eclare une urgence de $0{,}001$.

Trois \'epreuves ind\'ependantes en ont d\'elimit\'e la port\'ee~: elles ont localis\'e un
d\'efaut pr\'ecis de la couche de pr\'evision, \'etabli que l'outil rep\`ere celui qui se
m\'eprend sur son \'etat et non celui qui le tait, et mesur\'e o\`u la simulation apporte
r\'eellement sa valeur.

\medskip

\noindent\textbf{Mots-cl\'es~:} jumeau num\'erique de cha\^ine d'approvisionnement,
ad\'equation d'urgence, biais cognitif, aide multicrit\`ere \`a la d\'ecision, calibration de
pr\'evisions, biais d'automatisation, \emph{serious game}.
\end{otherlanguage}

\vspace{0.3cm}

\section*{Abstract}
\label{sec:resume_en}

In a supply chain, the urgency an operator declares often diverges from the urgency computable
from the real operational state. Over-declaring wastes shared capacity; under-declaring leaves
rupture risk invisible. No digital twin, published or patented, measures that gap yet.

This mission, carried out at the Direction de l'Innovation of ALTEN, designed and tested a
software layer that measures it, delivered as \textbf{SupplyScore}. Real urgency is
aggregated from six indicator blocks by a
\emph{non-compensatory} rule: the most degraded block drives the score toward its maximum, and
no healthy block can offset it. It is then propagated across the network. Declared urgency is
elicited separately, through a consistency-gated pairwise questionnaire, and the gap between
the two is scored on $[0,100]$.

In the serious-game tests carried out, over the six turns before a crisis becomes public, the
gap flags three nodes, all three of which later miss a delivery commitment, with no false alarm
on any turn; the worst-scoring of them is flagged eight turns before it misses, while its
manager declares an urgency of $0.001$.

Three independent tests then established its domain: they located a precise defect in the
forecast layer, established that the tool detects a declarant who misjudges their own
state rather than one who knows it and stays silent, and measured where the simulation
genuinely adds value.

\medskip

\noindent\textbf{Keywords:} supply chain digital twin, urgency adequacy, cognitive bias,
multi-criteria decision analysis, forecast calibration, automation bias, serious game.

\vspace{0.3cm}

\section*{Resumen}
\label{sec:resume_es}

\begin{otherlanguage}{spanish}
En una cadena de suministro, la urgencia que declara un operador difiere a menudo de la
urgencia calculable a partir del estado operativo real. Declarar de m\'as desperdicia
capacidad compartida; declarar de menos deja invisible el riesgo de ruptura. Ning\'un gemelo
digital, publicado o patentado, mide todav\'ia esa diferencia.

Esta misi\'on, realizada en la Direcci\'on de Innovaci\'on de ALTEN, dise\~n\'o y prob\'o una
capa de software que la mide, entregada con el nombre de \textbf{SupplyScore}. La urgencia
real se agrega a partir de seis bloques de
indicadores mediante una regla \emph{no compensatoria}: el bloque m\'as degradado lleva la
puntuaci\'on hacia su m\'aximo y ning\'un bloque sano puede compensarlo. Despu\'es se propaga
por la red. La urgencia declarada se obtiene por separado, mediante un cuestionario de
comparaciones por pares con control de coherencia, y la diferencia entre ambas se punt\'ua en
$[0,100]$.

En las pruebas de \emph{serious game} realizadas, en los seis turnos previos a que una crisis
se haga p\'ublica, la diferencia se\~nala tres nodos, los tres futuros incumplidores, sin una
sola falsa alarma; el peor puntuado se se\~nala ocho turnos antes de incumplir su hito,
mientras su responsable declara una urgencia de $0{,}001$.

Tres pruebas independientes delimitaron su alcance: localizaron un defecto preciso en la capa
de predicci\'on, establecieron que la herramienta detecta a quien se equivoca sobre su propio
estado y no a quien lo conoce y calla, y midieron d\'onde aporta valor realmente la
simulaci\'on.

\medskip

\noindent\textbf{Palabras clave:} gemelo digital de cadena de suministro, adecuaci\'on de
urgencia, sesgo cognitivo, an\'alisis multicriterio de decisi\'on, calibraci\'on de
predicciones, sesgo de automatizaci\'on, \emph{serious game}.
\end{otherlanguage}
